%═══════════════════════════════════════════════════════════════════════════════
\section{Introduction}\label{sec:introduction}
%═══════════════════════════════════════════════════════════════════════════════

\subsection{From Prediction Sets to Scores}\label{sec:intro_sets}

Consider a decision-maker who must act before an uncertain outcome is known.
They know which outcomes may occur, which actions are available to take, and what each
action would cost in each outcome. What they do not know is which outcome will
occur, so they ask a statistician for evidence.

Call whatever the statistician sends back a report. It might be a shortlist of
the outcomes still worth taking seriously, or a score for each outcome saying
how strongly the data count against it. Whatever its form, a report must be
valid, meaning it must not mislead more often than a pre-agreed error rate allows.

The established way to build such a report is to build a
prediction set, a shortlist of the outcomes that remain plausible. Conformal
prediction constructs such sets with a coverage guarantee: for a chosen level
$\alpha\in(0,1)$, the shortlist $S$ contains the realised outcome $Z$ with
probability at least $1-\alpha$, that is, $P[Z\in S]\geq1-\alpha$
\citep{vovk2005,angelopoulos2023conformal}. A cautious decision-maker can then
protect against the worst thing the shortlist still permits,
\begin{equation}
  \delta_S\in\argmin_{d\in\mathcal D}\ \max_{z\in S}L(d,z),
  \label{eq:intro_set_rule}
\end{equation}
writing $L(d,z)$ for what action $d$ costs when outcome $z$ occurs. Whenever
the realised outcome lies in $S$, the loss of $\delta_S$ cannot exceed the
worst loss over $S$, so coverage carries over to the decision,
\begin{equation}
  P\Big[L(\delta_S,Z)\leq\max_{z\in S}L(\delta_S,z)\Big]\geq1-\alpha .
  \label{eq:intro_set_guarantee}
\end{equation}
This is the basis of certified-decision methods
\citep{AndrewsChen2025,Kiyani2025}.

This guarantee is genuine, but a shortlist records only membership. Once two
outcomes are both in it, the set no longer says which one the data support
more strongly, so the rule treats a barely plausible outcome and a highly
likely one identically. The recommended action can also change with the
coverage level even though nothing about the decision has changed.

E-values keep what a shortlist discards. Rather than sorting each outcome in
or out, they give every outcome its own score $\mathcal E(z)\geq0$, large
where the data argue against that outcome and small where they do not. The
collection of these scores, one per outcome, is what the statistician hands
over. This thesis calls it a profile and writes it $\mathcal E$.

What stops a profile from claiming strong evidence everywhere is a single
requirement,
\begin{equation}
  \mathbb E_{P}[\mathcal E(Z)]\leq 1 ,
  \label{eq:intro_validity}
\end{equation}
that the scores average to at most one under the distribution $P$ the evidence
is weighed against. Section~\ref{sec:decision} extends this to a whole set
$\mathcal P$ of distributions.

A profile satisfying this restriction is valid, and a valid profile is also
called an e-variable. The condition acts as a budget: raising the score on one outcome forces a reduction elsewhere. In return, the scores can be compared against any threshold, even one chosen after they are seen, without losing the guarantee (Section~\ref{sec:evar}). This is the basis of safe
testing \citep{Shafer2021,Grunwald2024safe,Ramdas2023}. Fuzzy prediction sets
carry this graded evidence back to the prediction problem, replacing the
in-or-out verdict of a shortlist with a score for every outcome
\citep{KoningvanMeer2025,Vovk2025,Gauthier2025}.

\subsection{Two Ways to Judge a Profile}\label{sec:intro_two}

Validity rules out overclaiming, but it leaves a great deal open. Many
profiles satisfy it, and they need not lead the decision-maker to the same
action. The statistician must still choose one, and could let what each
mistake would cost guide that choice. This is where the existing literature
and the decision-maker part company.

The standard construction compares two distributions. One is $P$, the
distribution validity is required to hold under. The other is the belief $\widehat P_1$,
what the decision-maker expects to happen. Their ratio
\begin{equation}
  r(z):=\frac{\widehat P_1(z)}{P(z)}
  \label{eq:intro_lr}
\end{equation}
is large where the belief rates an outcome likelier than $P$ does, and the profile is built by passing this ratio through an increasing
shape function $w$, reporting $\mathcal E(z)=w(r(z))$. 

The usual measure of how much a profile tells us is its expected
log-evidence under the belief, $\mathbb E_{\widehat P_1}[\log\mathcal E(Z)]$,
which is high when large scores sit where the belief puts its probability.
Whenever $P$ gives every outcome positive probability, $w(r)=r$ maximises it
among valid profiles \citep{Grunwald2024safe,Koning2025continuous}. Statisticians with other preferences over how evidence accumulates choose
other shape functions, and \citet{KoningvanMeer2025} determine which $w$ is
best once such a preference is fixed.

However, all of this concerns how much the evidence tells us. It says nothing about
what happens when the evidence must produce an action. For that,
\citet{KoningvanMeer2025} score each action by
\begin{equation}
  C(d;\mathcal E):=\max_{z\in\mathcal Z}\frac{L(d,z)}{\mathcal E(z)} ,
  \label{eq:intro_cert}
\end{equation}
its largest loss across outcomes, with each loss first divided by the evidence
against the outcome that would produce it. Validity turns this single number
into a bound on the action's expected loss, and the decision-maker takes
whichever action has the smallest one.

Dividing by the score is how evidence lowers the bound. A large score on an
outcome shrinks the loss the action would suffer there, while a small score
leaves that loss large. Because of the budget, the scores cannot all be large,
so a profile has to choose where its large scores go.
The bound is the worst loss after division, so it stays low only if the large
scores sit where the action's losses are large.

This is where the standard construction can go wrong. It puts its large
scores on the outcomes the belief expects, those that $\widehat P_1$ rates
likelier than $P$ does, and is therefore left with small scores on the
outcomes the belief considers unlikely. If the action is expensive in
one of those unlikely outcomes, that loss is not reduced, and the action's
bound is poor. Which outcomes the belief expects and where the
decision-maker has most to lose are separate questions, and there is no
reason their answers should coincide.

An evidence profile can thus be judged in two ways, by the information it preserves, or
by the protection it gives to a decision. The two judgements need not agree, which leads to
the research question of the thesis:
\begin{quote}
Given a loss function and a stated attitude toward risk, how should a
statistician choose the e-variable supplied to a decision-maker?
\end{quote}

\subsection{Contribution}\label{sec:intro_contrib}

Existing theory chooses a profile by a preference over evidence, and cannot
say which preference represents a given decision-maker's loss. The closest predecessor is
the e-posterior of \citet{Grunwald2023epost}, which supports frequentist
decision-making without a prior and already divides a loss by an e-value in
the manner of \eqref{eq:intro_cert}. Neither supplies a design step in which the loss
function is an input to the construction of the profile itself. This step is the contribution of
this thesis, which uses three results to answer the question.

The first asks how low the bound on a chosen action's loss can be pushed. The
answer is to make the scores proportional to what that action costs in each
outcome, and the bound then equals that action's largest expected loss across
the distributions validity must hold under (Theorem~\ref{thm:cert}). No other
valid profile does as well, and the optimum is unique. Call it the safe
profile. It is built from the loss and those distributions alone, while the
belief plays no part in it.

The second result separates informing from inducing. The safe profile
always leads to the action it protects, unless some other action is cheaper in
every outcome, but the most informative profile, the likelihood ratio, need
not, even when the target is the action the belief
itself favours. Keeping as much information as possible while still reaching
the target therefore requires imposing the target as a constraint, and the
optimum then sits where the target ties with a competitor
(Proposition~\ref{prop:boundary}).

The third asks whether the standard construction could have found the
safe profile anyway. Since it reads each outcome only through $r$, two
outcomes with the same likelihood ratio must receive the same score, however
differently they cost. The safe profile is built from the loss
and can separate them. The two agree exactly when the loss is determined by
the ratio (Proposition~\ref{prop:no_shape}). Otherwise, no shape function
recovers the safe profile, so the loss must enter the design of the
profile rather than be applied to its output.